\item \textbf{¿Bajo qué condiciones se puede migrar un atributo de algún tipo de entidad que participa en un tipo de relación binaria y convertirse en un atributo del tipo de relación? ¿Cuál sería en el efecto?}

Las condiciones bajo las que se puede migrar un atributo de una relación binaria son cuando tenemos atributos cuyo valor depende de la combinación de las dos entidades relacionadas y no de únicamente una de ellas, es decir, cuando tenemos un atributo que describe a la instancia de la unión.

El efecto de esto es que el atributo deja de estar asociado directamente con una instancia de la entidad y pasa a estar asociada con cada instancia de la relación.
